% ============================================================
%  STILE DEL CHEAT SHEET DI INGEGNERIA DEL SOFTWARE
% ============================================================
% Si carica dopo preamble/igsw.tex, da cui prende colori, stili TikZ dei diagrammi UML (\umlclass, classe1, aggr,
% compos, gen, realizza, navig, dep, molt, notauml...) e l'aspetto dei listati. Aggiunge:
%   \argomento{Titolo}{nota}[titolo breve]  sezione numerata (1, 2, ...); il titolo breve va nel piè di pagina
%   \voce[larghezza del disegno]{etichetta}{Nome}{nota}{disegno}{significato}
%                       una voce di notazione: a sinistra numero e nome, al centro il disegno UML, a destra il
%                       significato (e, quando c'è, l'equivalente Java)
%   ambiente caso[opzioni]{etichetta}{Titolo}
%                       un caso «dal codice al diagramma»: a sinistra il codice Java (un lstlisting con
%                       style=javacs), dopo \tcblower il diagramma e la \regola{...}
%   ambiente struttura[opzioni]{etichetta}{Titolo}
%                       come caso, ma con il diagramma a sinistra (la struttura di un pattern) e il testo a destra
%   \scheda{Titolo}{righe}  la carta d'identità di un pattern: righe \pintento, \ptesto, \pcodice, \pdiagramma,
%                       \pvarianti, \pconfondere
%   ambiente guida{Titolo}  riquadro con una procedura o una tabella per scegliere
%   \vv{etichetta}       rimando cliccabile a una voce o a un caso: scrive il suo numero
%   \oltre               etichetta «oltre le slide» per ciò che nelle slide non c'è
% Voci e casi condividono la numerazione dell'argomento (3.1, 3.2...); \vv porta alla voce citata.

% ------------------------------------------------------------
% Colori
% ------------------------------------------------------------
\colorlet{ncolore}{swteal!85!black}    % voci di notazione
\colorlet{nsfondo}{swteal!5}           % sfondo alternato delle voci
\colorlet{ccolore}{swblue!85!black}    % casi «dal codice al diagramma»
\colorlet{pcolore}{ph3!85!black}       % pattern
\colorlet{acolore}{sworange!85!black}  % «attenzione», «da non confondere», «oltre le slide»
\colorlet{casocolore}{ccolore}         % colore dei riquadri caso e struttura: nella parte II diventa pcolore

% Pagine di apertura delle parti come nel formulario di MSPI: con titlesec la parte non lascia dietro di sé una
% pagina bianca (la classe è a due facciate, e il \part standard ne aggiungerebbe una)
\titleformat{\part}[display]
  {\centering\color{primary}}
  {\fontsize{20}{24}\selectfont\Gobold Parte \thepart}
  {18pt}
  {\color{primary!80!black}\fontsize{32}{40}\selectfont\Gobold #1}
  [\vspace{10pt}{\color{primary!40}\rule{0.35\linewidth}{2pt}}]
\assignpagestyle{\part}{empty}

% ------------------------------------------------------------
% Argomenti e numerazione
% ------------------------------------------------------------
\newcounter{argomento}
\newcounter{voce}
\renewcommand{\thevoce}{\arabic{argomento}.\arabic{voce}}
% Nome dell'ancora del PDF: diverso per ogni voce (voce.2.4, voce.3.4...), altrimenti i collegamenti alle voci di
% argomenti diversi con lo stesso numero finirebbero tutti sulla prima
\renewcommand{\theHvoce}{\arabic{argomento}.\arabic{voce}}

% \argomento{Titolo}{nota sotto il titolo}[titolo breve per il piè di pagina]
% Ogni argomento comincia su una pagina nuova: si ritrova subito sfogliando, e il titolo non resta mai in fondo a
% una pagina staccato dal primo riquadro (i riquadri non si spezzano, e quello che segue il titolo può essere alto).
\NewDocumentCommand{\argomento}{m m o}{%
  \clearpage
  \stepcounter{argomento}\setcounter{voce}{0}%
  \section{\arabic{argomento}\enspace #1}%
  \IfNoValueTF{#3}{\gdef\currentchapter{#1}}{\gdef\currentchapter{#3}}%
  \nobreak\vspace{-4mm}{\Lato\small\color{swgray}#2}\par\nobreak\vspace{2mm}%
}

% Numero in un rettangolino colorato
\newcommand{\badge}[2][ncolore]{%
  {\setlength{\fboxsep}{1.5pt}\colorbox{#1}{\Lato\bfseries\footnotesize\color{white}\strut\,#2\,}}}

% Rimando cliccabile a una voce o a un caso
\newcommand{\vv}[1]{\hyperref[v:#1]{\textcolor{swblue!80!black}{\Lato\bfseries\ref*{v:#1}}}}

% Etichetta per ciò che nelle slide non c'è (casi aggiunti per essere pronti alle novità)
\newcommand{\oltre}{%
  {\setlength{\fboxsep}{0.8pt}\fcolorbox{acolore}{white}{\Lato\upshape\mdseries\scriptsize\color{acolore}\,oltre le slide\,}}}

% Codice nel testo e testo dei diagrammi scritto come nel diagramma
\newcommand{\uml}[1]{{\FiraCodeNL\small #1}}

% ------------------------------------------------------------
% Le voci di notazione
% ------------------------------------------------------------
\newtcolorbox{vocebox}[1][]{%
  enhanced, breakable=false, boxrule=0pt, frame hidden, arc=0pt, colback=white,
  left=3mm, right=3mm, top=2mm, bottom=2mm, before skip=0pt, after skip=0pt,
  borderline south={0.5pt}{0pt}{swgray!30},
  before upper={\setlength{\parskip}{0pt}},
  #1}

\newlength{\vocenome}
\setlength{\vocenome}{36mm}
% \voce[larghezza del disegno]{etichetta}{Nome}{nota sotto il nome}{disegno}{significato}
\NewDocumentCommand{\voce}{O{62mm} m m m +m +m}{%
  \ifodd\value{voce}\def\vocesfondo{nsfondo}\else\def\vocesfondo{white}\fi
  \begin{vocebox}[colback=\vocesfondo]
    \begin{minipage}[c]{\vocenome}\raggedright
      \refstepcounter{voce}\label{v:#2}%
      \badge[ncolore]{\thevoce}\par\vspace{2pt}%
      {\bfseries #3\par}%
      \ifx\relax#4\relax\else{\footnotesize\color{swgray}#4\par}\fi
    \end{minipage}\hfill
    \begin{minipage}[c]{#1}\centering #5\end{minipage}\hfill
    \begin{minipage}[c]{\dimexpr\linewidth-\vocenome-#1-8mm\relax}\raggedright\small\setlength{\parskip}{3pt}#6%
    \end{minipage}%
  \end{vocebox}}

% ------------------------------------------------------------
% I casi «dal codice al diagramma»
% ------------------------------------------------------------
% Il riquadro ha il titolo con il numero; sotto, due colonne: a sinistra il codice (un lstlisting, quindi
% l'ambiente non può essere il corpo di un comando), a destra, dopo \tcblower, il diagramma e la regola.
\newtcolorbox{casobox}[2][]{%
  enhanced, breakable=false, sidebyside, sidebyside align=center, sidebyside gap=4mm,
  lefthand width=96mm, colback=white, colframe=casocolore!40, boxrule=0.5pt, arc=1.5mm,
  left=2mm, right=3mm, top=2mm, bottom=2mm, before skip=2.5mm, after skip=2.5mm,
  colbacktitle=casocolore!8, coltitle=casocolore, fonttitle=\bfseries, toptitle=0.8mm, bottomtitle=0.8mm,
  titlerule=0.5pt, titlerule style={casocolore!30}, segmentation hidden,
  before upper={\setlength{\parskip}{0pt}}, before lower={\setlength{\parskip}{0pt}\centering},
  title={#2}, #1}
\NewDocumentEnvironment{caso}{O{} m m}{%
  \begin{casobox}[#1]{\refstepcounter{voce}\label{v:#2}\badge[casocolore]{\thevoce}\enspace #3}%
}{\end{casobox}}
% Un esempio completo di un pattern: come un caso, ma il programma si compila e si esegue (cheatsheet/
% verifica_listati.py confronta l'output con il listato «Terminale» a destra, sotto il diagramma)
\NewDocumentEnvironment{programma}{O{} m m}{%
  \begin{casobox}[#1]{\refstepcounter{voce}\label{v:#2}\badge[casocolore]{\thevoce}\enspace #3}%
}{\end{casobox}}
% La regola sotto il diagramma: come si passa dal codice al disegno
\newcommand{\regola}[1]{\par\vspace{2.5mm}{\raggedright\footnotesize\setlength{\parskip}{2pt}#1\par}}
% La struttura di un pattern: il diagramma a sinistra, i partecipanti a destra
\NewDocumentEnvironment{struttura}{O{} m m}{%
  \begin{casobox}[lefthand width=100mm, before upper={\setlength{\parskip}{0pt}\centering},
                  before lower={\setlength{\parskip}{3pt}\raggedright\small}, #1]%
    {\refstepcounter{voce}\label{v:#2}\badge[casocolore]{\thevoce}\enspace #3}%
}{\end{casobox}}

% Listati del cheat sheet: come quelli della dispensa, ma senza numeri di riga e più compatti. I nomi di metodi e
% tipi da colorare li raccoglie preamble/nomi_listati.py nel file cheatsheet/listati-nomi.tex.
% File generato da preamble/nomi_listati.py: non modificarlo a mano.
% Metodi (giallo) e tipi (verde acqua) che compaiono nei listati Java del cheat sheet.
\lstdefinelanguage{JavaVSCnomiCS}[]{JavaVSC}{
  morekeywords=[3]{
    add,aggiungi,allOf,annota,append,applicaSconto,area,avvia,avvisa,between,build,calcola,
    calcolaPrezzo,casella,cerca,chiudi,click,compra,compraNuova,conferma,deposita,descrivi,di,
    diAutore,distanza,equals,esegui,espressa,fattura,flick,gestisci,get,getArea,getCreati,getEta,
    getIndirizzo,getInstance,getMessage,getNome,getNumero,getOrDefault,getPrezzo,getProdotto,
    getQuantita,getResult,getStampante,getYears,grafico,incassa,includeOS,invia,iscrivi,isFedele,
    join,main,metti,monta,nasce,normale,now,of,ofNullable,pcGaming,pcUfficio,preleva,premi,prendi,
    print,println,prioritario,pubblica,put,registra,remove,ridimensiona,rimuovi,runAsync,saluta,
    salva,scala,setNome,setVeicolo,size,smonta,somma,sposa,sqrt,ssd,stampa,supera,toString,vendi,
    verificaPin,volume,withDedicatedGPU,withSSD},
  morekeywords=[4]{
    @Override,Aeroporto,Archivio,ArchivioFile,ArrayDeque,ArrayList,Articolo,Ascoltatore,
    Assemblaggio,Auto,Automobile,Badge,Biblioteca,Brano,BuilderPC,Calcolatrice,Casa,Casella,
    Cassiere,Cerchio,Cliente,Collection,CompletableFuture,Computer,ComputerBuilder,Configuratore,
    Configurazione,Contabilita,Contatore,Contatto,Conto,Corso,Counter,Cubo,Deque,Documento,Email,
    Esame,Esecutore,Exception,Figura,Finestra,Forma,Forma3D,Formato,GestioneOrdini,Giocatore,
    Grafico,HashMap,HashSet,IllegalArgumentException,IllegalStateException,Immagine,Impiegato,
    Indirizzo,Integer,Lampada,Libro,LinkedHashSet,LinkedList,List,LocalDate,Magazzino,Main,Map,Math,
    Moto,Negozio,Notificatore,Optional,Ordine,Partita,Passeggero,Period,Persona,Playlist,Pratica,
    Prenotazione,Preventivo,Prodotto,Pulsante,Punto,Queue,Registro,Report,Rettangolo,Ricetta,
    Ricevuta,Riga,Rubrica,Ruota,SaldoInsufficiente,Saluto,Salvabile,Scacchiera,Scatola,
    SchedaTecnica,Servizio,Set,SingletonClass,Sportello,Squadra,Stampabile,Stampante,Stanza,Stato,
    String,StringBuilder,Studente,System,Taglia,Torneo,TreeSet,Utente,Veicolo,Vendite,Void,Volo,
    Voto},
}

\lstdefinestyle{javacs}{
  style=java, language=JavaVSCnomiCS,
  numbers=none, xleftmargin=8pt, framexleftmargin=4pt, framesep=4pt, xrightmargin=4pt,
  aboveskip=0pt, belowskip=0pt,
}
% nell'output di un programma l'apostrofo è testo, non l'inizio di una stringa (Attendi l'erogazione)
\lstdefinestyle{shellcs}{
  style=shell, deletestring={[b]'}, numbers=none, xleftmargin=8pt, framexleftmargin=4pt, framesep=4pt, xrightmargin=4pt,
  aboveskip=0pt, belowskip=0pt,
}

% ------------------------------------------------------------
% La carta d'identità di un pattern
% ------------------------------------------------------------
\newtcolorbox{schedabox}[1]{%
  enhanced, breakable=false, colback=white, colframe=pcolore!45, boxrule=0.5pt, arc=1.5mm,
  left=3mm, right=3mm, top=1.5mm, bottom=2.5mm, before skip=2.5mm, after skip=2.5mm,
  colbacktitle=pcolore!9, coltitle=pcolore, fonttitle=\bfseries, toptitle=0.9mm, bottomtitle=0.9mm,
  titlerule=0.5pt, titlerule style={pcolore!30}, title={#1},
  before upper={\setlength{\parskip}{0pt}},
}
\newcommand{\scheda}[2]{\begin{schedabox}{#1}#2\end{schedabox}}

% Una riga della scheda: etichetta a sinistra, testo a destra
\newlength{\schedaetichetta}
\setlength{\schedaetichetta}{30mm}
\newcommand{\rigascheda}[3]{%
  \par\vspace{2.5pt}\noindent
  % \noindent prima del colore: il riferimento delle minipage [t] deve essere la prima riga di testo
  \begin{minipage}[t]{\schedaetichetta}\noindent{\Lato\bfseries\small\color{#1}#2}\end{minipage}%
  \begin{minipage}[t]{\dimexpr\linewidth-\schedaetichetta\relax}\setlength{\parskip}{3pt}\raggedright\noindent#3%
  \end{minipage}\par}
\newcommand{\pintento}[1]{\rigascheda{pcolore}{Intento}{#1}}
\newcommand{\ptesto}[1]{\rigascheda{pcolore}{Nel testo}{#1}}
\newcommand{\pcodice}[1]{\rigascheda{pcolore}{Nel codice}{#1}}
\newcommand{\pdiagramma}[1]{\rigascheda{pcolore}{Nel diagramma}{#1}}
\newcommand{\pvarianti}[1]{\rigascheda{ccolore}{Varianti}{#1}}
\newcommand{\pconfondere}[1]{\rigascheda{acolore}{Da non confondere}{#1}}

% ------------------------------------------------------------
% Riquadri «guida»: procedure e tabelle per scegliere
% ------------------------------------------------------------
\newtcolorbox{guidabox}[2][]{%
  enhanced, breakable=false, colback=primary!4, colframe=primary!55, boxrule=0.5pt, arc=1.5mm,
  left=3mm, right=3mm, top=2mm, bottom=2.5mm, before skip=2mm, after skip=3mm,
  fonttitle=\bfseries, colbacktitle=primary!55, coltitle=white, title={#2},
  before upper={\setlength{\parskip}{4pt}}, #1}
\NewDocumentEnvironment{guida}{O{} m}{\begin{guidabox}[#1]{#2}}{\end{guidabox}}
% Riquadro «attenzione»: un errore tipico
\newtcolorbox{attenzione}[1]{%
  enhanced, breakable=false, colback=acolore!5, colframe=acolore!60, boxrule=0.5pt, arc=1.5mm,
  left=3mm, right=3mm, top=1.5mm, bottom=2mm, before skip=2mm, after skip=3mm,
  fonttitle=\bfseries, colbacktitle=acolore!12, coltitle=acolore, title={#1},
  before upper={\setlength{\parskip}{3pt}}}

% ------------------------------------------------------------
% Piccoli elementi dei diagrammi usati spesso nel cheat sheet
% ------------------------------------------------------------
\tikzset{
  % associazione semplice (linea continua) e testo dei ruoli
  assocl/.style={line width=0.7pt, draw=black!80},
  ruolo/.style={font=\tiny, inner sep=1.5pt, text=black!85},
  % linea tratteggiata verso una nota o una classe associativa
  tratt/.style={dashed, line width=0.6pt, draw=black!70},
  % contenimento di una classe annidata: cerchio con la croce dal lato della classe esterna
  annida/.style={{Circle[open, length=2.6mm, width=2.6mm]}-, line width=0.7pt, draw=black!80},
}
% Croce del lato non navigabile: \noNav{(punto)}{angolo della linea}
\newcommand{\noNav}[2]{%
  \draw[line width=0.75pt, black!80] ($#1+(#2+45:0.13)$) -- ($#1+(#2+225:0.13)$)
                                     ($#1+(#2+135:0.13)$) -- ($#1+(#2+315:0.13)$);}
% Croce dentro il cerchio del contenimento: \piuAnnida{(centro del cerchio)}
\newcommand{\piuAnnida}[1]{%
  \draw[line width=0.6pt, black!80] ($#1+(-0.13,0)$) -- ($#1+(0.13,0)$) ($#1+(0,-0.13)$) -- ($#1+(0,0.13)$);}
% Stereotipo sopra il nome di una classe: \stereotipo{interface}{Nome}
\newcommand{\stereotipo}[2]{{\tiny\normalfont «#1»}\\#2}
% Molteplicità e ruolo all'estremo di un'associazione, appena fuori dalla classe e lontano dalla punta della
% freccia. Il primo argomento è il punto in cui la linea tocca la classe; la lettera dice da dove arriva la linea:
%   \estO da sinistra (ovest), \estE da destra (est), \estN dall'alto (nord), \estS dal basso (sud).
% La molteplicità va sopra (o a destra), il ruolo sotto (o a sinistra).
\newcommand{\estO}[3]{%
  \node[molt, anchor=south east] at ($#1+(-0.1,0.07)$) {#2};
  \node[ruolo, anchor=north east] at ($#1+(-0.1,-0.07)$) {#3};}
\newcommand{\estE}[3]{%
  \node[molt, anchor=south west] at ($#1+(0.1,0.07)$) {#2};
  \node[ruolo, anchor=north west] at ($#1+(0.1,-0.07)$) {#3};}
\newcommand{\estN}[3]{%
  \node[molt, anchor=south west] at ($#1+(0.08,0.1)$) {#2};
  \node[ruolo, anchor=south east] at ($#1+(-0.08,0.1)$) {#3};}
\newcommand{\estS}[3]{%
  \node[molt, anchor=north west] at ($#1+(0.08,-0.1)$) {#2};
  \node[ruolo, anchor=north east] at ($#1+(-0.08,-0.1)$) {#3};}
